% Chapter 6: Digit sums and the alternative Kummer formulation

\section{Digit extraction}

\begin{definition}[Base-$p$ digit]
  \label{def:digit}
  \rocq{MathCompKummer.digit_sum.digit}
  For $p \geq 2$, the $i$-th digit of $n$ in base $p$ is
  \[
    d_i(n) := \left\lfloor \frac{n}{p^i} \right\rfloor \bmod p.
  \]
\end{definition}

\begin{lemma}[Digit bound]
  \label{lem:digit_bound}
  \uses{def:digit, lem:ltn_mod}
  \rocq{MathCompKummer.digit_sum.digit_bound}
  For $p \geq 2$,
  \[
    d_i(n) < p.
  \]
\end{lemma}

\begin{lemma}[Digit representation]
  \label{lem:digits_repr}
  \uses{def:digit, lem:divn_eq}
  \rocq{MathCompKummer.digit_sum.digits_repr}
  For $p \geq 2$ and $n > 0$,
  \[
    n = \sum_{i=0}^{\texttt{trunc\_log}(p,n)} d_i(n) \cdot p^i.
  \]
\end{lemma}

\section{Digit sum}

\begin{definition}[Digit sum]
  \label{def:digit_sum}
  \uses{def:digit}
  \rocq{MathCompKummer.digit_sum.digit_sum}
  The sum of base-$p$ digits of $n$:
  \[
    \digitsum_p(n) := \sum_{i=0}^{\texttt{trunc\_log}(p,n)} d_i(n).
  \]
\end{definition}

\section{Legendre via digit sums}

\begin{lemma}[Legendre's formula via digit sums]
  \label{lem:legendre_digit_sum}
  \uses{lem:logn_fact, def:digit_sum}
  \rocq{MathCompKummer.digit_sum.legendre_digit_sum}
  For prime $p$ and $n > 0$,
  \[
    (p-1) \cdot v_p(n!) = n - \digitsum_p(n).
  \]
  \begin{proof}
    From the digit representation $n = \sum_i d_i(n) \cdot p^i$
    and Legendre's formula $v_p(n!) = \sum_{k \geq 1} \lfloor n/p^k \rfloor$.
    A telescoping argument gives the result.
  \end{proof}
\end{lemma}

\section{Kummer via digit sums}

\begin{theorem}[Kummer's theorem — digit sum form]
  \label{thm:kummer_digit_sum}
  \uses{thm:kummer, lem:legendre_digit_sum}
  \rocq{MathCompKummer.digit_sum.kummer_digit_sum}
  For prime $p$ and $0 < k \leq n$,
  \[
    (p-1) \cdot v_p\!\binom{n}{k}
    = \digitsum_p(k) + \digitsum_p(n-k) - \digitsum_p(n).
  \]
  \begin{proof}
    Apply Lemma~\ref{lem:legendre_digit_sum} to $n!$, $k!$, and $(n-k)!$,
    then use $v_p\binom{n}{k} = v_p(n!) - v_p(k!) - v_p((n-k)!)$.
  \end{proof}
\end{theorem}
